\begin{textsample}{NEG Dim 4 – 2016\_c5 – Score 1.00 – 2016\_c5/enade\_2016\_biomed\_7425d9844e.txt}  \label{ex:f4_neg_253}
CAPÍTULO 5 – DISTRIBUIÇÃO DOS CONCEITOS 5.1 - PANORAMA NACIONAL DA DISTRIBUIÇÃO DOS CONCEITOS O Conceito Enade é um indicador de qualidade que avalia o desempenho dos estudantes a partir dos resultados do Enade.
Os valores possíveis vão de 1 ( pior situação ) a 5 ( melhor situação ).
As Unidades de Observação com apenas um ou sem nenhum concluinte participante não obtêm o Conceito Enade, ficando Sem Conceito ( SC ).
A Tabela 5.1 apresenta a quantidade e distribuição de cursos de Biomedicina participantes do Enade/2016, por faixa de conceito e Grande Região.
Enfatiza -se, mais uma vez, que a diferença entre os cursos tabulados neste capítulo e no capítulo 2 corresponde aos cursos sem conceito, ou seja àqueles, em princípio, sem alunos concluintes que houvessem participado da prova ou àqueles que tivessem tão somente um aluno concluinte19.
Observando -se os dados da Tabela 5.1, nota -se que, dos 235 cursos participantes, 108 ( 46,0 \% ) classificaram -se com conceito 3, o valor modal.
Esse foi também o conceito modal em três regiões : Sudeste ( 47,9 \% ), Sul ( 57,5 \% ) e Centro-Oeste ( 42,9 \% ).
Nas demais regiões, a moda foi o conceito 2, atribuído para 42,1 \% dos cursos da região Norte e para 50,0 \% dos cursos da região Nordeste.
O conceito 2 foi o segundo mais frequente em nível nacional ( 24,3 \%, correspondendo a 57 cursos ) e o conceito 4, o terceiro ( 16,6 \%, correspondendo a 39 cursos ).
Houve, ainda, 17 cursos ( 7,2 \% ) que receberam conceito 5 e 11 cursos ( 4,7 \% ) que receberam conceito 1.
Dos 235 cursos de Biomedicina, três ( 1,3 \% ) ficaram sem conceito ( SC ). 19 Estes cursos com somente um concluinte não foram considerados no capítulo 2.
A região Norte participou com 19 cursos ou 8,1 \% do total nacional.
Desses, oito cursos receberam o conceito 2, o que equivale a 42,1 \% do total regional, além de ser o conceito modal na região, como já comentado.
A quatro cursos ( 21,1 \% ) atribuiu -se conceito 1 e a sete cursos ( 36,8 \% ), conceito 3.
Nenhum curso ficou sem conceito ( SC ) ou conceito 4 ou 5 nessa região.
A região Nordeste participou com 38 cursos ou 16,2 \% do total nacional.
Nessa região, 19 cursos ( 50,0 \% em termos regionais ) obtiveram conceito 2, o conceito modal para a região, como já comentado.
Os conceitos 1 e 3 foram atribuídos, respectivamente, a um e a 13 cursos ( 2,6 \% e 34,2 \%, respectivamente ).
Ainda nessa região, três cursos ( 7,9 \% ) foram avaliados com o conceito 4, e dois cursos ( 5,3 \% ), com conceito 5.
Nenhum curso ficou sem conceito ( SC ).
Dos 117 cursos participantes da região Sudeste ( 49,8 \% do total nacional ), o conceito 3 foi o conceito modal, como já anteriormente assinalado, tendo sido obtido por 56 cursos ( 47,9 \% ).
O conceito 1 foi atribuído a quatro cursos ( 3,4 \% ) e o conceito 2, a 22 cursos ( 18,8 \% ).
Receberam os conceitos 4 e 5, respectivamente, 23 cursos ( 19,7 \% ) e nove cursos ( 7,7 \% ).
Nessa região, três cursos ( 2,6 \% ) ficaram sem conceito ( SC ).
Os 40 cursos da região Sul corresponderam a 17,0 \% do total nacional.
Como já indicado, houve predominância do conceito 3, atribuído a 57,5 \% dos cursos da região ( conceito modal ), ou seja, atribuído a 23 dos 40 cursos participantes na região Sul.
O conceito 2 foi atribuído a dois cursos ( 5,0 \% ), o conceito 4, a nove cursos ( 22,5 \% ) e o conceito 5, a seis cursos ( 15,0 \% ).
Nenhum curso recebeu conceito 1 ou ficou sem conceito ( SC ) na região Sul.
Já dos 21 cursos participantes na região Centro-Oeste ( 8,9 \% do total nacional ), nove receberam conceito 3, o conceito modal ( 42,9 \% ), como já fora destacado.
A dois cursos ( 9,5 \% ) foi atribuído conceito 1, a seis cursos ( 28,6 \% ), conceito 2 e quatro cursos, conceito 4.
Nessa região, nenhum curso ficou sem conceito ( SC ) ou com conceito 5.
O Gráfico 5.1 apresenta a distribuição do Conceito Enade segundo Grande Região : poligonais mais à esquerda representam Grandes Regiões com pior distribuição, e poligonais mais à direita, Grandes Regiões com melhores conceitos.
A região Norte ( linha azul ) apresenta os cursos com conceitos mais baixos, ao passo que a região Sul ( linha roxa ) apresenta a distribuição com valores maiores.
Cursos em IES situadas nas regiões Nordeste ( linha verde ), Sudeste ( linha vermelha ) e Centro-Oeste ( linha preta ) apresentam uma situação intermediária. 5.2 - CONCEITOS POR CATEGORIA ADMINISTRATIVA E GRANDE REGIÃO A Tabela 5.2 apresenta a distribuição dos cursos participantes do Enade/2016 de Biomedicina, por Categoria Administrativa, de acordo com os conceitos por eles alcançados, segundo as Grandes Regiões brasileiras.
Dos 235 cursos participantes, 26 ( 11,1 \% ) eram ministrados em Instituições Públicas e 209 ( 88,9 \% ), em Instituições Privadas.
De acordo com as informações da Tabela 5.2, em termos nacionais, entre Instituições Públicas e Privadas, observa -se que, dos 17 cursos avaliados com conceito 5, sete eram oferecidos em IES Públicas e dez, em IES Privadas.
Dos 26 cursos participantes de IES Públicas, o conceito 4 foi o valor modal, atribuído a dez cursos ( 38,5 \% ).
Entre os demais cursos de Instituições Públicas participantes, um curso ( 3,8 \% ) foi avaliado com conceito 2, oito cursos ( 30,8 \% ) obtiveram conceito 3 e, como já comentado, sete cursos ( 26,9 \% ) receberam conceito 5.
Nesta categoria, nenhum curso ficou sem conceito ( SC ) ou com conceito 1.
Na rede privada, o conceito modal foi 3, atribuído a 100 cursos, o correspondente a 47,8 \% dos 209 cursos da categoria.
Entre os demais cursos participantes, 11 ( 5,3 \% ) receberam conceito 1, e 56 ( 26,8 \% ), conceito 2.
O conceito 4 foi atribuído a 29 cursos ( 13,9 \% ) e o conceito 5, a dez cursos ( 4,8 \% ).
Nesta Categoria Administrativa, três cursos ( 1,4 \% ) ficaram sem conceito ( SC ).
Na análise por região, observa -se que, na região Norte, as Instituições Privadas participaram com 18 cursos ( 94,7 \% do total regional ), a proporção foi mais elevada do que a registrada nas demais regiões brasileiras.
Dos quais nenhum ficou sem conceito ( SC ) ou com conceito 4 ou 5.
O conceito modal para as Instituições Privadas na região foi 2, atribuído a oito cursos, correspondendo a 44,4 \% dessa categoria na região.
Quatro cursos foram avaliados com conceito 1 e seis cursos, com conceito 3.
As Instituições Públicas participaram com somente um curso na região Norte ( 5,3 \% do total regional ), este curso recebeu conceito 3.
Na região Nordeste, a rede privada concentrou 34 dos 38 cursos participantes, o equivalente a 89,5 \% do total da região.
Dentre os cursos oferecidos por IES Privadas no Nordeste, 19 cursos foram avaliados com conceito 2, o conceito modal.
Três cursos obtiveram conceitos 1, 4 e 5 ( um curso, cada ) e 12 cursos receberam conceito 3.
Nessa combinação de categoria e região, nenhum curso ficou sem conceito ( SC ).
As Instituições Públicas da região Nordeste participaram com quatro cursos ( 10,5 \% do total da região ).
Desses, dois obtiveram conceito 4, o valor modal, seguidos por outros dois cursos, que receberam conceitos 3 e 5 ( um curso, cada ).
Nenhum curso ficou sem conceito ( SC ), e nenhum curso de IES Pública recebeu conceito 1 ou 2 nessa região.
Na região Sudeste, a proporção de cursos da rede privada foi 90,6 \%, correspondendo a 106 dos 117 cursos participantes.
Nessa categoria e região, o conceito modal foi 3, atribuído a 53 cursos.
Três cursos não receberam conceito ( SC ).
Os demais foram avaliados com conceito 1 ( quatro cursos ), conceito 2 ( 21 cursos ), conceito 4 ( 18 cursos ) e conceito 5 ( sete cursos ).
Entre os 11 cursos oferecidos por Instituições Públicas na região Sudeste, o conceito modal foi 4, atribuído a cinco cursos.
Os demais cursos receberam os conceitos 2 ( um curso ), 3 ( três cursos ) e 5 ( dois cursos ).
Na região Sudeste, nesta categoria, nenhum curso ficou sem conceito ( SC ) ou recebeu conceito 1.
As Instituições Privadas concentraram 33 dos 40 cursos participantes da região Sul, 82,5 \% do total regional, a menor das proporções desta categoria entre as regiões.
Desses, 22 obtiveram conceito 3, o conceito modal.
Nesta combinação de Categoria Administrativa e Grande Região, os conceitos 2 e 4 foram atribuídos a dois e a sete cursos, respectivamente.
Dois cursos receberam conceito 5, e nenhum recebeu conceito 1 ou ficou sem conceito ( SC ).
As Instituições Públicas na região Sul participaram com sete cursos, 17,5 \% dos cursos da região, a maior das proporções desta categoria nas regiões.
O conceito modal foi 5, atribuído a quatro cursos.
Aos demais cursos foram atribuídos os conceitos 3 ( um curso ) e 4 ( dois cursos ).
Nenhum curso da região Sul oferecido por IES Pública ficou sem conceito ( SC ) ou recebeu os conceitos 1 e 2.
Na região Centro-Oeste, 18 dos 21 cursos participantes eram de Instituições Privadas ( 85,7 \% em termos regionais ).
Desses, sete receberam conceito 3, o conceito modal.
Entre os demais, dois receberam conceito 1, seis receberam conceito 2 e três receberam conceito 4.
Nenhum curso ficou sem conceito ( SC ) ou com conceito 5.
Quanto aos três cursos oferecidos por Instituições Públicas na região Centro-Oeste ( 14,3 \% do total regional ), nota -se que dois cursos ficaram com conceito 3, conceito modal, e um curso recebeu conceito 4.
Nenhum curso ficou sem conceito ( SC ) ou com os demais conceitos. receberam conceito 4.
Nenhum curso ficou sem conceito ( SC ) ou com conceito 5.
Quanto aos três cursos oferecidos por Instituições Públicas na região Centro-Oeste ( 14,3 \% do total regional ), nota -se que dois cursos ficaram com conceito 3, conceito modal, e um curso recebeu conceito 4.
Nenhum curso ficou sem conceito ( SC ) ou com os demais conceitos.
O Gráfico 5.2 apresenta a distribuição do Conceito Enade segundo a Categoria Administrativa da IES.
Os cursos em IES Públicas ( linha azul ) apresentam uma poligonal mais à direita do que os em IES Privadas ( linha verde ) e, consequentemente, uma distribuição de Conceitos Enade com valores maiores. 5.3 - CONCEITOS POR ORGANIZAÇÃO ACADÊMICA E GRANDE REGIÃO Na Tabela 5.3, encontra -se a distribuição dos conceitos atribuídos aos cursos participantes do Enade/2016 na Área de Biomedicina, por Organização Acadêmica, segundo as Grandes Regiões brasileiras.
Dos 235 cursos de Biomedicina participantes, 97 eram oferecidos em Universidades, 50 em Centros Universitários e 88 em Faculdades.
Esta distribuição corresponde a, respectivamente, 41,3 \%, 21,3 \% e 37,4 \% dos cursos.
De acordo com os dados apresentados, de todos os 17 cursos avaliados com conceito 5, 13 eram vinculados a Universidades.
Os cursos vinculados a esse tipo de Organização Acadêmica tiveram o conceito 3 como conceito modal, atribuído a 42 cursos ( 43,3 \% ).
Os demais cursos vinculados a Universidades receberam os conceitos 1 ( um curso ), 2 ( 13 cursos ), 4 ( 27 cursos ) e conceito 5 ( 13 cursos, como já mencionado ).
Ficou sem conceito ( SC ), somente um curso.
Entre os cursos em Centros Universitários, o conceito modal foi também 3, atribuído a 31 cursos ( 62,0 \% ).
Um curso vinculado a esse tipo de Organização Acadêmica ficaram sem conceito ( SC ), e os outros cursos receberam os conceitos 1 ( um curso ), 2 ( seis cursos ), 4 ( nove cursos ) e 5 ( dois cursos ).
Dos 88 cursos mantidos por Faculdades, 38 ( 43,2 \% ) receberam conceito 2, o conceito modal.
O segundo conceito mais frequentemente atribuído aos cursos oferecidos por essa categoria de IES foi o conceito 3, recebido por 35 cursos.
Na sequência, nove cursos obtiveram conceito 1 e três cursos, conceito 4.
O conceito 5 foi atribuído a dois cursos, e um curso ficou sem conceito ( SC ).
Considerando -se separadamente as regiões brasileiras, verifica -se que, na região Norte, as Universidades concentraram dois dos 19 cursos participantes, no qual um ficou com conceito 2 e um com conceito 3.
Os Centros Universitários da região Norte foram representados por três cursos, aos quais foram todos atribuídos ao conceito 3.
As Faculdades participaram com 14 cursos na região Norte que receberam os conceitos 1 ( quatro cursos ), 2 ( sete cursos, conceito modal ) e 3 ( três cursos ).
Nenhuma Faculdade da região ficou sem conceito ( SC ) ou obteve conceito 4 ou 5.
Na região Nordeste, as Universidades participaram com oito dos 38 cursos da Área de Biomedicina oferecidos na região, dos quais nenhum ficou sem conceito ( SC ) ou com conceito 1.
O conceito modal foi 3, atribuído a três cursos, enquanto o conceito 5 foi atribuído a um curso.
Quatro cursos obtiveram os conceitos 2 e 4 ( dois cursos, cada ).
Os Centros Universitários contaram com sete cursos participante na região Nordeste.
Esses obtiveram conceitos 2 ( dois cursos ) e 3 ( quatro cursos, o conceito modal ).
As Faculdades foram representadas por 23 cursos na região Nordeste, 15 dos quais obtiveram o conceito modal 2.
O conceito 1 foi atribuído a um curso e o conceito 4, também, a um curso.
Seis cursos obtiveram o conceito 3.
Na região Nordeste, nenhum curso oferecido por Faculdades recebeu conceito 5, e nenhum curso ficou sem conceito ( SC ).
Na região Sudeste, as Universidades concentraram 61 dos 117 cursos de Biomedicina da região.
Entre os cursos em Universidades, nessa região, o conceito modal foi 3, atribuído a 26 cursos.
Os demais cursos receberam os conceitos 1 ( um curso ), 2 ( dez cursos ), 4 ( 15 cursos ) e 5 ( oito cursos ).
Um curso ficou sem conceito.
Os Centros Universitários participaram com 26 cursos na região Sudeste, dos quais 14 obtiveram o conceito modal 3.
O conceito 1 foi atribuído a um único curso, o conceito 2, a quatro cursos e o conceito 4, a seis cursos.
Um curso ficou sem conceito ( SC ) e nenhum curso ficou com conceito 5 nessa combinação de organização acadêmica e região.
Já as Faculdades foram representadas por 30 cursos na região Sudeste, um dos quais ficou sem conceito ( SC ).
Aos demais cursos foram atribuídos os conceitos 1 ( dois cursos ), 2 ( oito cursos ), 3 ( 16 cursos, o conceito modal ), 4 ( dois cursos ) e 5 ( um curso ).
Dos 40 cursos da região Sul, 19 eram de Universidades, para os quais o conceito modal foi 3, atribuído a oito cursos.
Nesse tipo de organização, nenhum curso ficou sem conceito ( SC ) ou recebeu conceito 1 ou 2.
Os demais cursos receberam os conceitos 4 ( sete cursos ) e 5 ( quatro cursos ).
Dos dez cursos participantes de Centros Universitários da região Sul, nenhum ficou sem conceito ( SC ) ou recebeu conceito 1 ou 2.
O conceito modal, 3, foi atribuído a sete cursos.
Os conceitos restantes, 4 e 5, foram atribuídos, respectivamente, a dois e a um curso.
Quanto aos 11 cursos vinculados a Faculdades na região Sul, oito receberam o conceito modal 3.
O conceito 2 foi atribuído a dois cursos, e o conceito 5, a um curso.
Nenhum curso recebeu conceito 1 ou 4, nem tampouco ficou sem conceito ( SC ).
Na região Centro-Oeste, sete dos 21 cursos eram de Universidades, dos quais quatro receberam conceito 3, conceito modal, e três ficaram com conceito 4.
Os Centros Universitários da região Centro-Oeste contaram com quatro cursos, sendo três \textbf{alocados} ao conceito 3, conceito modal, e um ao conceito 4.
Dos dez cursos oferecidos em Faculdades na região Centro-Oeste, seis receberam o conceito modal 2.
Aos demais cursos foram atribuídos os conceitos 1 ( dois cursos ) e 3 ( dois cursos ).
Nenhuma Faculdade da região Centro-Oeste ficou sem conceito ( SC ) ou com conceito 4 ou 5.
O Gráfico 5.3 apresenta a distribuição do Conceito Enade segundo a Organização Acadêmica da IES.
Os cursos em Universidades ( linha azul ) apresentam uma poligonal mais à direita do que as demais, denotando uma melhor distribuição de Conceitos Enade.
No outro extremo, temos os cursos oferecidos em Faculdades ( linha vermelha ) com a poligonal mais à esquerda e, consequentemente, uma distribuição de Conceitos Enade com menores valores.
IES de Centros Universitários ( linha verde ) apresentam uma situação intermediária.

% matched lemmas: alocar
\end{textsample}
